\section{Implementation and its checking boundary}
\label{sec:implementation}
The implementation separates a numerical producer from replay. The producer uses nonlinear and MILP solvers to seek useful evidence. Replay reconstructs the exact model, checks the cuts, matches the rational master, and checks its discrete derivations. Reusing the producer's search history is unnecessary. Several mathematical utilities are shared, so replay is independent of the numerical search but is not a separately developed implementation of every mathematical rule.

\begin{figure}[tb]
\centering
\begin{tikzpicture}[
  box/.style={draw,rounded corners,align=center,text width=3.0cm,
    minimum height=1.8cm,inner sep=4pt,font=\small},
  node distance=0.45cm,>={Stealth}]
\node[box] (search) {Numerical search\\OA points and slopes\\Exact MILP proof producer};
\node[box,right=of search] (bundle) {Saved bundle\\Model and cut lemmas\\Rational master\\Completed proof};
\node[box,right=of bundle] (check) {Complete replay\\Exact model and cuts\\Master identity\\Every proof inference};
\node[box,right=of check] (bound) {Bound for the\\original loaded model\\$s f_{\rm orig}(x)\geq\beta$\\for every $x\in F$};
\draw[->] (search) -- (bundle);
\draw[->] (bundle) -- (check);
\draw[->] (check) -- (bound);
\end{tikzpicture}
\caption{The certificate interface. Complete replay requires every checking obligation to pass; the last implication is objective-preserving bound transfer. Search is outside the trusted base once its output passes. Lean verifies the mathematical contract and typed reference checkers separately, with no verified connection to the Python implementation or benchmark certificates.}
\label{fig:architecture}
\end{figure}

\subsection{Artifact and execution contract}
\label{sec:artifact-contract}
A certificate bundle consists of the supplied model, a nonlinear-lemma file, a rational LP master, and a completed VIPR proof. Each nonlinear lemma names its normalized row and supplies a rational point $z$, a sparse rational slope vector $a$, and an intercept $b$. The checker reconstructs the box and expressions from the model; it does not trust a supplied box, derivative value, curvature label, or reported solver bound. New bundles record the semantic identifier \texttt{pyomo-expression-rationals-v1} and source hash. For compatibility with the saved historical files these fields may be absent; a present identifier must match, and a present source hash is checked. In either case the explicit model passed to replay defines the theorem being checked. The frozen campaign manifest additionally identifies its exact bytes.

Complete checking performs the following operations in order.
\begin{enumerate}
\item Load and validate the model, normalize the objective and nonlinear rows, check domains and curvature, extract exact affine rows, and propagate rational bounds.
\item Reconstruct each lemma's row and symbolic derivatives, verify its point and variable names, and recompute a rigorous intercept bound. Require the supplied $b$ to be no larger than that bound.
\item Regenerate the LP master and require byte equality with the saved LP file. Compare the VIPR problem against the reconstructed rational master semantically, as specified below.
\item Check every supplied master solution and every discrete proof step. Require a finite minimization lower bound for the normalized master. If an external checker was requested, require matching successful corroboration as well.
\end{enumerate}
Only then does the public routine return \texttt{ok=True} and a bound in the original objective sense. Partial mode returns \texttt{status="partial"}, potentially with \texttt{partial\_ok=True}, but \texttt{ok=False} and no certified bound. A missing proof, exception, unsupported input, or failed check cannot expose a certified numerical result. The lower-level VIPR kernel supports infeasibility and either objective sense; the MINLP interface deliberately supports finite normalized lower-bound reports. Theorems about infeasibility in Section~\ref{sec:soundness} do not enlarge this API.

The model loader executes a Python/Pyomo file and applies a documented dedenting repair to top-level model assignments in converted files. It therefore belongs to trusted input construction. Replay requires one active objective and refuses active disjunctive, logical, or special ordered set (SOS) constraints. Supported source variable and affine-row names follow the identifier pattern of a letter or underscore followed by letters, digits, or underscores; reserved auxiliary names cannot alias original variables. The exact extractor replaces fixed variables by their exact values and performs polynomial coefficient accumulation with rational arithmetic. Nonpolynomial subtrees remain available for original-domain checks, even when multiplied by zero. Differentiation and interval evaluation use the same rational-leaf interpretation, rather than a separate floating-point coefficient aggregation.

\subsection{Sufficient curvature and domain rules}
\label{sec:curvature-rules}
Curvature checking combines composition rules in the spirit of disciplined convex programming \citep{grant2006-dcp} with a small number of algebraic recognizers. Each expression carries a curvature classification and, where justified, rational lower and upper range bounds. These are separate facts: a product may be nonnegative without being convex or concave. Unknown curvature or a missing range endpoint is not replaced by a guessed value.

For completeness, the composition principle is immediate from Jensen's inequality. If $v$ is convex and $u$ is convex and nondecreasing on the range of $v$, apply monotonicity to $v(\lambda x+(1-\lambda)y)\leq\lambda v(x)+(1-\lambda)v(y)$ and then convexity of $u$. If $u$ is convex and nonincreasing, use a concave $v$ instead. Reversing both convexity inequalities gives the corresponding concave rules. The implementation applies these principles to sums, rational scalings, affine absolute values, and the functions in Table~\ref{tab:composition}. Negative scaling interchanges convex and concave classifications.

\begin{table}[htbp]
\centering\small
\caption{Principal scalar composition tests. ``Positive'' means a strictly positive certified lower bound; all original expression domains must also pass.}
\label{tab:composition}
\begin{tabular}{p{0.19\linewidth}p{0.55\linewidth}p{0.14\linewidth}}
\toprule
Expression & Sufficient condition & Curvature\\
\midrule
$\exp(v)$ & $v$ affine or convex & Convex\\
$\log(v)$ & $v$ positive, affine or concave & Concave\\
$\sqrt v$ & $v$ nonnegative, affine or concave & Concave\\
$v^p$, $p\geq1$ & $v$ nonnegative, affine or convex & Convex\\
$v^p$, $0<p<1$ & $v$ nonnegative, affine or concave & Concave\\
$v^p$, $p<0$ & $v$ positive, affine or concave & Convex\\
$v^{2k}$, $k\in\Z_{>0}$ & $v$ affine, or nonnegative convex, or nonpositive concave & Convex\\
$k_0/v$, $k_0>0$ & $v$ positive, affine or concave & Convex\\
$k_0^v$, $k_0>0$ & $v$ affine; also $v$ convex if $k_0>1$, or concave if $k_0<1$ & Convex\\
\bottomrule
\end{tabular}
\end{table}
The power and reciprocal tests follow from first and second derivatives on the stated intervals; for $k_0^v$ use $\exp(v\log k_0)$. Power one preserves curvature, and power zero is the constant one after checking the original base domain. Nonnegative-domain rules extend to an included zero endpoint only when the expression is defined and continuous there. Convexity recognition at such an endpoint does not establish differentiability there.

\paragraph{Quadratic forms and norms.}
For $q(x)=x^\top Qx+b^\top x+c$ with rational symmetric $Q$, global convexity on $\R^n$ is equivalent to $Q\succeq0$ because $\nabla^2q=2Q$. This gives a sufficient test on any certified box, including a box with fixed coordinates. Off-diagonal monomial coefficients are split equally between $Q_{ij}$ and $Q_{ji}$. The exact elimination test uses the following elementary characterization.
\begin{lemma}[Exact quadratic tests]
\label{lem:quadratic-recognition}
For symmetric $Q$, a positive leading pivot $\alpha$ in
$Q=\left[\begin{smallmatrix}\alpha&v^\top\\v&C\end{smallmatrix}\right]$
reduces positive semidefiniteness to $C-vv^\top/\alpha\succeq0$. A negative diagonal entry rules it out. A zero diagonal entry requires its whole remaining row and column to be zero, in which case they can be removed. These rules give an exact rational test. Moreover, if
\[
 H=\begin{bmatrix}Q&b/2\\b^\top/2&c\end{bmatrix}\succeq0,
\]
then $\sqrt{q(x)}$ is globally defined and convex.
\end{lemma}
\begin{proof}
For $\alpha>0$, completing the square writes the quadratic form as
$\alpha(t+v^\top z/\alpha)^2+z^\top(C-vv^\top/\alpha)z$.
A negative diagonal gives a negative value on a coordinate vector. If a diagonal is zero and some corresponding entry $v_j$ is nonzero, the restriction to those two coordinates has a term $2v_jt$ with no $t^2$ term and takes negative values. This proves the elimination rules by induction. For the second statement, factor the real positive semidefinite matrix as $H=R^\top R$. Then $\sqrt{q(x)}=\|R(x,1)\|_2$, the norm of an affine map, and the triangle inequality proves convexity. The factor need not be rational or computed by the checker; exact testing of $H$ is sufficient.
\end{proof}
Testing $-Q$ gives concavity. The homogenized test is used only as a sufficient recognition rule for square roots of quadratics; no claim is made to recognize every convex square-root expression. At a norm origin, finite subgradients may exist even though the symbolic-gradient route abstains.

\paragraph{Monomials and geometric means.}
The algebraic recognizer collects rational exponents in
$m(x)=\prod_i x_i^{\alpha_i}$ on a positive orthant. It accepts convexity when all $\alpha_i\leq0$, or when exactly one exponent is positive and $\sum_i\alpha_i\geq1$. It accepts concavity when all $\alpha_i\geq0$ and $\sum_i\alpha_i\leq1$. A rational coefficient scales or reverses the result. These established conditions are stated by \citet[Theorems~2.1--2.2, p.~507]{lundell2009-signomial}, who attribute them to earlier work. We include a direct proof to make the recognition contract self-contained and avoid reliance on floating eigenvalue calculations.
\begin{lemma}[Monomial tests]
\label{lem:monomial-recognition}
The preceding sufficient curvature tests hold on the positive orthant. When all nonzero exponents are positive, the resulting convexity or concavity extends continuously to the nonnegative orthant wherever the product is so defined.
\end{lemma}
\begin{proof}
With $D=\operatorname{diag}(1/x_i)$, differentiation gives
\[
 \nabla^2m(x)=m(x)D\bigl(\alpha\alpha^\top-\operatorname{diag}\alpha\bigr)D.
\]
If every $\alpha_i\leq0$, the middle matrix is the sum of a positive semidefinite outer product and a nonnegative diagonal matrix. For nonnegative exponents, weighted Cauchy--Schwarz yields
$(\sum_i\alpha_i u_i)^2\leq(\sum_i\alpha_i)\sum_i\alpha_i u_i^2$,
so the middle matrix is negative semidefinite if the exponent sum is at most one.

For the remaining convex case remove zero exponents and write $\alpha=(p,-b_1,\ldots,-b_k)$ with $b_i>0$, $B=\sum_i b_i$, and $p\geq1+B$. The lower block of the middle matrix is $\operatorname{diag}(b)+bb^\top\succ0$, and
$b^\top(\operatorname{diag}(b)+bb^\top)^{-1}b=B/(1+B)$,
as follows by multiplying that matrix by $\mathbf 1/(1+B)$. Its Schur complement is $p(p-1)-p^2B/(1+B)=p(p-1-B)/(1+B)\geq0$. If $k=0$, this is simply $p(p-1)\geq0$. The Hessian criteria prove the interior assertions. Taking limits of the convexity inequalities proves their continuous boundary extensions.
\end{proof}
The syntax recognizer restricts fractional powers of products to justified nonnegative variable factors and a unit constant before distributing the exponents. Original denominator and power domains are checked first. Thus it does not use identities such as $\sqrt{x^2}=x$ on a negative domain. For example, the product-range calculation can prove $xy\geq0$ for $x,y\geq0$, after which the exponent test recognizes the concavity of $\sqrt{xy}$. It never labels the bilinear product itself concave for that reason.

\paragraph{One-variable linear fractions.}
A further rule recognizes $(\alpha x+\beta)/(\gamma x+\delta)$ when the denominator is bounded strictly away from zero with a constant sign. For $\gamma\ne0$, write
\[
 \frac{\alpha x+\beta}{\gamma x+\delta}
 =\frac{\alpha}{\gamma}+\frac{k}{\gamma x+\delta},\qquad
 k=\beta-\frac{\alpha\delta}{\gamma}.
\]
Its second derivative is $2k\gamma^2/(\gamma x+\delta)^3$. Convexity therefore holds when $k$ and the denominator have the same sign, concavity when their signs differ, and the function is constant if $k=0$. A constant denominator uses rational scaling instead. This is a one-variable test; it is not a general curvature test for ratios or perspectives. The former perspective shortcut was disabled because curvature of the inner expression on its own variable bounds need not justify curvature on the actual ratio domain.

\paragraph{Domain and range enforcement.}
Range propagation uses rational endpoint arithmetic for sums, scales, integer powers, and finite products. Finite products use all four endpoint products. On sign-certified half-lines the corresponding extreme products give the finite endpoint, and unsupported sign combinations return an unknown range. Exponential, logarithmic, and rational-power endpoint bounds use outward interval arithmetic at 128-bit precision, with exact rational extraction of their dyadic endpoints. Precision affects how often a strict sign can be proved, not the logical sign criterion.

The original tree is visited before simplification. Division and negative integer powers require the relevant range to exclude zero strictly; logarithms require positivity. Noninteger constant powers require a nonnegative base for positive exponents and a positive base for negative exponents. Variable powers require a positive base. A square root requires nonnegativity, possibly supplied by Lemma~\ref{lem:quadratic-recognition}. Only the stated arithmetic nodes and supported unary functions are admitted. These checks take place on declared bounds, not on a domain inferred from nonlinear constraints. A model can therefore be mathematically convex and feasible yet be rejected by the sufficient rules.

\subsection{Interval cut replay and numerical proposals}
\label{sec:cut-implementation}
The nonlinear expression is translated to SymPy with exact rational leaves, and symbolic derivatives are built from that same expression. The interval evaluator implements addition, multiplication, powers, exponential, logarithm, and absolute value. A rational point is converted by interval division of its integer numerator and denominator, never through a binary64 approximation. Rational powers are evaluated with the applicable square-root or exponential/logarithm formula. If a formula or derivative is unsupported, undefined, or nonfinite at the point, the cut is refused. In particular, recognizing an absolute value or norm as convex does not provide a general nonsmooth derivative rule.

The precise bridge from differentiation to \eqref{eq:g-support} is a derivative along feasible segments, not merely a list of finite coordinate differences. If, for each $x\in B$, the right derivative of $\theta\mapsto g(z+\theta(x-z))$ at zero is $p^\top(x-z)$, convexity gives
\[
 g(x)-g(z)\geq \lim_{\theta\downarrow0}
 \frac{g(z+\theta(x-z))-g(z)}{\theta}=p^\top(x-z).
\]
A differentiable extension near $z$ is sufficient; valid relative-domain chain rules can also give this condition at a boundary. The symbolic-expression layer is trusted to provide derivatives with this meaning. Finiteness of arbitrary one-sided coordinate derivatives alone would not suffice: $-\sqrt{xy}$ on the nonnegative quadrant has zero derivatives along the two axes at the origin, but the zero vector does not support it. The implemented symbolic quotient derivatives are singular there and the cut is rejected. Likewise, singular derivatives of $-\sqrt{x}$ and unsupported sign expressions at norm origins are refused. Removable singularities can be simplified only consistently with the original checked domain. These are explicit obligations of the expression and differentiation implementation, not consequences of interval finiteness or a mechanized verification of every accepted expression.

Cut checking uses a 200-bit interval context. All function and derivative endpoints must be finite. The right-hand side of \eqref{eq:safe-intercept} is enclosed by interval operations, including the maximum of its two endpoint products on a finite coordinate. The final endpoint's sign, integer mantissa, and binary exponent are read directly and converted to a rational. This avoids rounding the endpoint through a different precision before comparing it to the supplied intercept. The arithmetic libraries and symbolic derivative implementation remain trusted components of this operation.

The numerical producer uses a separate loaded model. Its OA log supplies candidate row names, points, and slopes; exact reconstructed rows define the certificate mathematics. Points are rationalized and clipped to the propagated box. Slopes are initially shortened to 12 significant decimal digits by default, with directed repairs on one-sided coordinates and exact slopes where available on free coordinates. The intercept is rounded downward to 30 significant decimal digits after rigorous evaluation. These digit counts control representation and practical replay robustness; they provide no universal separation margin between two independent interval evaluations. Complete replay still requires the exact inequality $b\leq$ the newly computed lower endpoint, with no tolerance. Poor points or broad enclosures can lead to fewer cuts or weaker bounds.

\subsection{Binding the discrete proof to the master}
\label{sec:master-identity}
The LP writer includes rational affine rows, accepted cut rows, explicit finite bounds, integrality declarations, and the normalized objective. Every free variable has explicit infinite bounds. Fixed variables retain equal bounds, and a nonzero affine objective constant is carried by a new variable \texttt{objconst} fixed to one. This prevents a silent loss of the constant during proof export. An epigraph coordinate is real and free. Integer bounds tightened by propagation are justified by Lemma~\ref{lem:propagation}.

The VIPR problem parser is shared with the proof kernel. Comparison requires the same variable set under a bijection, the same integer-variable set, minimization sense, and identical rational objective coefficients. A uniform single \texttt{t\_} prefix introduced by SCIP can be removed only if this yields that bijection; genuine original names with that prefix are preserved when already matched. Constraints and bounds are compared as multisets after dividing each nonzero row by the absolute value of its first coefficient in sorted variable order. This permits positive rational row scaling, retaining the relation direction. The current matcher is deliberately stricter than general linear equivalence: it does not recognize arbitrary equivalent row transformations or negative scaling even of equalities. Additional, missing, or altered rows are rejected. A certificate for a different presolved model cannot pass merely because it has the same reported optimum.

The saved LP must also equal the regenerated LP byte for byte. This adds an artifact consistency check; the semantic VIPR comparison supplies the mathematical link to the proof. Original zero-coefficient affine tautologies are omitted in both constructions, and contradictory constant rows are rejected earlier. A hash or LP byte comparison alone cannot replace proof-input matching.

\subsection{Proof grammar, arithmetic, and streaming}
\label{sec:proof-implementation}
The kernel accepts complete proofs in a restricted subset of VIPR versions 1.0 and 1.1. It requires the section order
\[
 \texttt{VER, VAR, INT, OBJ, CON, RTP, SOL, DER}
\]
with the declared numbers of variables, indices, rows, solutions, and derivations. Input uses ASCII integer or fraction tokens with positive denominators; decimal and exponent notation are outside this parser's contract. Variable names, integer indices, vector coefficient indices, and linear-combination references must satisfy the corresponding uniqueness and range checks, including entries whose value is zero. Omitted solution coordinates mean zero and are still checked against all original rows and integrality conditions. The finite primal endpoint of a range report requires a feasible solution attaining that endpoint or a better one.

Each noncomment derivation occupies one line and carries a complete \texttt{asm}, \texttt{sol}, \texttt{lin}, \texttt{rnd}, or \texttt{uns} reason implementing Proposition~\ref{prop:vipr-invariant}. Row labels are metadata; references identify rows by indices. The \texttt{OBJ} coefficient abbreviation denotes the exact parsed objective. Arithmetic uses rational numbers through python-flint. A combination with incompatible inequality directions or a stronger-than-justified right-hand side is rejected even if its discrepancy is numerically small. Assumptions arise only from actual assumption rows and are tracked through every nonzero inference term.

The checker requires every inference reference to precede its row, including references with zero multipliers. A lifetime annotation is $-1$ or a nonnegative in-range index; references after a stated last use are errors. The optional single terminal \texttt{global} marker emitted by SCIP is accepted only if the independently computed assumption set is empty. It cannot remove assumptions or validate a coefficient. Repeated markers, incomplete reasons, unexpected extra tokens, missing derivations, and noncomment suffix material beyond the declared count are rejected.

An assumption-free row dominating the requested dual endpoint or contradiction is recorded as the proving derivation. It need not be the final row. Every declared later derivation must still pass, and the parser must consume the entire file. Consequently a valid proof prefix followed by invalid arithmetic cannot be accepted. A later valid but weaker bound is harmless. The lower-level kernel can also accept a range with no finite dual endpoint when its remaining obligations hold; the MINLP interface requires a finite lower endpoint before reporting a result.

Large proofs are processed in two passes. A first pass validates record shapes, references, and lifetime annotations and computes actual last uses with two arrays of eight-byte integers, about 16 bytes per row. A second pass retains one such array and live sparse rows, freeing a row only after its computed final reference. Rows annotated $-1$ are therefore also freed when safe. The parser initially holds the original master rows for solution checking; this storage and the longest input line must be included in memory accounting. The method avoids storing the whole proof text, but has no constant-memory guarantee: many simultaneously live dense rows, long rational tokens, or a large original master can still consume substantial memory.

\subsection{Failures that motivated the checking boundary}
\label{sec:checker-repairs}
Two small counterexamples show why an external success message is insufficient. The inspected upstream \texttt{viprchk} checkout\footnote{VIPR checkout \texttt{30f2951d1e90e47afa821bdd1b12b82246656c42}; the audit records the source hash and regression inputs. The statement concerns those inspected bytes.} used an initially zero best objective for a solution inference even with an empty solution section. A matching artifact for $\min x$ subject to $x\geq1$ could thereby receive acceptance for a claimed lower bound of $10^6$, although $x=1$ is feasible. The new kernel refuses a solution cutoff without an exactly checked witness. Wood et al.\ also discuss solution-inference and reference-management weaknesses in older checker behavior \citep{wood2026-vipr-smt}; this work makes no priority claim for identifying that general risk.

The second reproduced defect rounded the continuous inequality $x\geq1/2$ to $x\geq1$. This implication is false, even when every printed coefficient is rational. The corrected rule checks both integer coefficients and integer variable declarations before applying a floor or ceiling. Likewise, complementary branch inequalities form an exhaustive disjunction only for an integral linear form. Merely seeing adjacent integer right-hand sides is insufficient.

Other targeted regressions cover exact floating-leaf aggregation, lost domains after cancellation, singular norm derivatives, omitted slopes on half-lines and free coordinates, objective constants and maximization signs, source/master mismatches, invalid solution witnesses, cross-branch assumption retention, duplicate entries, invalid lifetimes, forged \texttt{global} markers, and invalid suffixes after an already sufficient bound. These tests establish distinct observable contracts; passing them is not a formal proof of the parser or executable.

The distinction between rejection and refutation is material. An insufficient safe-cut enclosure can reject a globally valid cut. An unsupported curvature rule can reject a convex model. In contrast, an exact arithmetic audit can establish that a particular recorded inference is invalid. Even the latter does not prove the final bound false if some different valid proof exists. The experimental section therefore retains these outcomes separately rather than relaxing the check to recover historical success labels.

\subsection{Generation, frozen replay, and trust}
\label{sec:trust}
The current producer uses numerical OA with Gurobi and IPOPT, followed by exact SCIP and VIPR completion tools. It first tries its default exact-SCIP configuration and can retry with selected separation, presolve, and propagation features disabled. Each attempt has its own search budget, and proof completion and checking add time. A solver time limit is consequently not a total pipeline wall-clock limit. Fraction canonicalization in newly generated completed proofs preserves their rational values; solution sections and inference meanings are not rewritten to obtain acceptance. Existing historical artifact directories are preserved.

Replay requires only the checker dependencies: Pyomo, NumPy, SymPy, mpmath, and python-flint, together with the included Python modules. Neither Gurobi, IPOPT, nor SCIP is required to check an already complete bundle. The optional external \texttt{viprchk} must return a zero status and one complete range message matching the exact internal result; it is additional corroboration. The generation harness requests this corroboration, while standalone replay need not.

The campaign wrapper hashes model, lemma, LP, and proof files before and after checking. It records source and environment identifiers, input records, optional external-checker identity, explicit missing/rejected/verified/error/timeout outcomes, and the original record index. Resume requires matching artifact fingerprints and the inputs, source/environment identifiers, external checker, and timeout pinned in the manifest; worker parallelism may change. A worker runs in its own process group so its timeout also stops descendant checkers. Completed output records are flushed; only an interrupted final JSON fragment may be recovered. These measures provide provenance and detect ordinary changes; they do not establish semantic equivalence or protect against an adversary that changes and restores a file during verification. Stable artifact contents and trusted model execution remain explicit execution assumptions. Summaries report whether all expected record indices are present exactly once and use exact rational objective-sense comparisons before converting values for display. An incomplete cohort can still be summarized, with its completeness flag false.

The trusted computing base comprises the loader and exact reconstruction, domain/range and curvature rules, symbolic differentiation, interval arithmetic and endpoint conversion, rational propagation, serialization and master comparison, proof parser and rational kernel, their libraries, and the runtime. The numerical search and proof producer need not be trusted when their outputs pass this contract. Shared trusted code can create common failures, so separate replay and adversarial tests reduce dependence on the numerical solver without eliminating implementation trust. Focused proof-assistant theorems in the following section address only their stated mathematical coverage; they do not automatically certify this entire software stack or every benchmark artifact.
